\section{Related Work}
\label{sec:related}

\subsection{Underwater Obstacle Perception and Avoidance}

Reactive avoidance on physical underwater vehicles is established across
both sensing modalities, and the modality largely determines the failure
mode the method must absorb. Acoustic pipelines exploit direct range: recent
forward-looking-sonar systems combine a learned detector with temporal
association and line-of-sight guidance, and are validated in lake
trials~\cite{li2026underwat}, while sonar-based planners have long served as
the reference solution for turbid water~\cite{mane2024eroasdef}. Optical
pipelines instead recover geometry indirectly and must cope with the
resulting scale ambiguity: monocular optic-flow divergence supports
autonomous braking and safe-distance keeping on a real ROV across a range of
visibility conditions~\cite{bergantin2024realtime}; stereo geometry supports
a committed semicircular bypass around located obstacles with return to a
predefined line in a pool~\cite{zhang2022binocula}; omnidirectional visual
SLAM supports potential-field evasion on a real
AUV~\cite{ochoa2022collisio}; and learned end-to-end policies map monocular
imagery, optionally fused with a single-beam echo sounder, directly to
motion commands on BlueROV2-class vehicles in
pools~\cite{yang2022monocula,lin2024uivnavun,yang2025duvindif}. The common
limitation for our purposes is not capability but evidence structure: each
of these systems demonstrates that its own controller works on its own
vehicle, so the reported numbers characterize a method, not the simulator
that might have predicted it. We therefore adopt a deliberately conventional
optical avoidance stack, because the object under study here is the
simulator rather than the avoidance algorithm.

\subsection{Temporal Estimation and Risk-Aware Avoidance}

Filtering detections over time, and letting the resulting uncertainty shape
the safety decision, is likewise a mature idea. In multi-object tracking the
measurement-noise covariance is routinely made responsive to detector
output: a linear confidence scaling of the measurement noise originated in
NSA-Kalman~\cite{du2021giaotrac} and was popularized by
StrongSORT~\cite{du2023strongso}, while later work derives the covariance
from the empirical spread of candidate boxes~\cite{solanocarrillo2024utrackmu}
or from a learned regression head whose output distribution is scored for
quality~\cite{lee2024uncertai}. Underwater, the same instinct appears as
confidence-driven adaptation of tracker parameters to suppress acoustic
false positives~\cite{li2026underwat}, as an extended Kalman filter on an
optic-flow-derived relative distance~\cite{bergantin2024realtime}, and as
covariance-inflated collision-risk zones feeding a potential-field
planner~\cite{li2024anobstac}. On the control side, uncertainty is
propagated into the safety constraint itself: belief control barrier
functions bound collision risk under a sampled environment
belief~\cite{han2025riskaware,han2026riskawar}, and scenario-based model
predictive control bounds collision probability jointly over a planned
trajectory~\cite{groot2025scenario}. Two consequences follow for this paper.
First, neither ``we filter detections with a Kalman filter'' nor ``our
planner accounts for uncertainty'' can be offered as a contribution. Second,
what remains genuinely under-reported is the \emph{comparative} behavior of
these estimators under the specific disturbances a closed-loop avoidance run
produces --- detector dropout at maneuver onset, abrupt ego-motion, and
young-track outliers --- which is what our ablation isolates.

\subsection{Underwater Simulation and Sim-to-Real Transfer}

Marine simulators have matured to the point where the same autonomy stack can
address a simulated and a physical vehicle without modification. HoloOcean
provides vehicle dynamics, optical and acoustic sensor models, and a ROS~2
interface~\cite{potokar2022holoocea,potokar2024holoocea,romrell2025apreview},
and has been used in software- and hardware-in-the-loop configurations in
which a real vehicle computer runs against the simulator and the resulting
behavior is compared against field trials of the same
vehicle~\cite{meyers2025testinga}. Transfer results follow the same pattern
elsewhere: a high-fidelity vehicle model supports reinforcement-learning
avoidance that transfers to pool experiments with a reported 94\% simulated
success rate~\cite{zhang2025simtoreal}; a photogrammetric digital twin
supervises a policy on a physical BlueROV2 in a harbor, with a dynamic-window
planner as the comparator~\cite{mari2026deeprein,mari2025digitalt}; and
instrumented facilities pair a pool with a digital
twin~\cite{torroba2026marinari,buche2026lotusimm}. Complementary work
calibrates the simulator itself rather than the policy, by identifying
BlueROV2 hydrodynamics against measured
responses~\cite{benzon2022anopenso,orjales2025towardsa}. What these studies
report is transfer or model agreement: the trajectory matched, the controller
survived, the coefficients improved. What they do not report is whether the
simulator, used as an evaluation instrument, would have ordered competing
avoidance methods the way the water does. It is also worth noting how the
physical evidence is constituted: in the closest BlueROV2 study the obstacles
avoided during the sea trials exist only inside the digital twin, and
obstacle sensing is computed from the twin rather than from onboard
perception~\cite{mari2026deeprein}.

\subsection{Simulator Predictivity and the Remaining Gap}

The distinction we need is between successful transfer and predictive
simulation, and it has a precise formulation outside marine robotics. Kadian
et al.~\cite{kadian2020simrealp} run identical navigation code in a simulator
and in a scanned replica of the same physical space, and quantify predictivity
by the correlation between simulated and real per-method outcomes, using the
Sim-vs-Real Correlation Coefficient to ask whether the \emph{ranking} of
methods is preserved; tuning simulator parameters raised that correlation
from 0.18 to 0.844, which is the strongest available evidence that
predictivity is a property one can deliberately engineer. Related work
outside navigation ranks sensor-model fidelity by its effect on a downstream
estimator judged against replicated real
runs~\cite{mahajan2024quantify}, and validates camera models by the response
of a downstream detector rather than at pixel
level~\cite{elmquist2022aperform}. The counterweight is
Truong et al.~\cite{truong2022rethinki}: higher simulation fidelity produced
\emph{worse} transfer for learned navigation policies, so added fidelity is
an empirical question, not a monotone improvement.

Taken together, prior work establishes underwater visual and acoustic
obstacle avoidance on physical vehicles, probabilistic temporal filtering of
obstacle state, uncertainty-aware collision avoidance, BlueROV2
experimentation with external ground truth, high-fidelity underwater
simulation, HoloOcean-based software- and hardware-in-the-loop testing,
underwater sim-to-real transfer, and simulator-predictivity analysis in
ground robotics. The question that remains is not whether an avoidance
controller transfers. It is whether progressively calibrating an underwater
simulator makes it quantitatively more predictive of physical safety and
performance, and of the relative ranking of obstacle-avoidance methods, on a
real ROV. Table~\ref{tab:related} places the most directly comparable studies
against the axes that question requires.

%==============================================================================
% Related-work positioning table.  Every cell is sourced from
% docs/LITERATURE_AND_NOVELTY.md (deep-read set) or from the closed-list
% verification recorded in paper/STATUS.md.  Status: READY.
% Full-width table*: column widths are fixed so the row text wraps instead of
% overflowing the text block.
%==============================================================================
\begin{table*}[t]
\centering
\caption{Positioning against the most directly comparable studies.
``Phys.\ obst.'' means the avoided obstacle physically existed in the water.
``Sim--real calibration'' means the simulator was deliberately adjusted to
measured reality. ``Predictivity / ranking'' means the study quantified how
well simulation predicts real outcomes or preserves the ordering of methods.}
\label{tab:related}
\tiny
\setlength{\tabcolsep}{2pt}
\renewcommand{\arraystretch}{1.25}
\begin{tabular}{@{}
  p{0.56in} p{0.54in} p{0.80in} p{0.66in} p{0.62in}
  p{0.66in} p{0.92in} p{0.72in} p{0.70in}@{}}
\toprule
\textbf{Work} & \textbf{Vehicle} & \textbf{Perception} &
\textbf{Temporal est.} & \textbf{Avoidance} & \textbf{Simulator} &
\textbf{Physical validation} & \textbf{Sim--real calibration} &
\textbf{Predictivity / ranking} \\
\midrule
Bergantin \cite{bergantin2024realtime}
  & BlueBee ROV & Mono.\ optic flow & EKF (rel.\ distance)
  & Reactive braking & --- & Pool/harbour/lake, phys.\ obst.\ & --- & --- \\
Yang \cite{yang2022monocula}
  & BlueROV2 & Mono.\ + echo sounder & ---
  & End-to-end RL & Unity & Pool, phys.\ obst.\ & Domain random.\ & --- \\
UIVNAV \cite{lin2024uivnavun}
  & BlueROV2 Hv.\ & Mono.\ RGB & ---
  & Imitation policy & OysterSim & Tank, phys.\ obst.\ & --- & --- \\
DUViN \cite{yang2025duvindif}
  & BlueROV2 Hv.\ & Mono.\ RGB, depth feat.\ & ---
  & Diffusion policy & Unity & Tank, phys.\ obst., mocap GT & --- & --- \\
Alinei-Poian\u{a} \cite{alineipoiana2024abluerov}
  & BlueROV2 & Mono.\ RGB (YOLO) & 15-state EKF
  & --- (mapping) & --- & Pool, overhead cam.\ inside filter & --- & --- \\
Li \cite{li2026underwat}
  & AUV & Forward-looking sonar & Confidence-adaptive
  & RA-LOS guidance & --- & Lake, phys.\ obst.\ & --- & --- \\
Li \cite{li2024anobstac}
  & AUV & Abstracted detections & KF + covariance risk
  & Improved APF & Numerical & --- (simulation only) & --- & --- \\
Han \cite{han2025riskaware}
  & Underwater (sim.) & --- & Sampled belief
  & Belief CBF & Numerical & --- (simulation only) & --- & --- \\
Han \cite{han2026riskawar}
  & BlueROV2 & 3D sonar & SMC-PHD (RFS)
  & BCBF-QP filter & PyBullet & Tank, phys.\ obst., mocap GT & --- & --- \\
Mari \cite{mari2026deeprein}
  & BlueROV2 & Twin-derived occup.\ grid & ---
  & PPO vs.\ DWA & Unity twin & Harbor, \emph{virtual} obst.\
  & Photogramm.\ twin & --- \\
Zhang \cite{zhang2025simtoreal}
  & Underwater robot & --- & ---
  & RL & High-fidelity model & Pool, phys.\ obst.\
  & Model / I-O alignment & --- \\
Meyers \cite{meyers2025testinga}
  & CougUV & Simulated sensors & Vehicle nav.\ filter
  & Waypoints & HoloOcean 2.0, SIL/HIL & Field trials
  & Single fidelity & Descriptive comparison \\
Kadian \cite{kadian2020simrealp}
  & Ground robot & RGB-D & ---
  & 9 learned policies & Habitat & Scanned real lab
  & Post-hoc param.\ tuning & \textbf{SRCC ranking} \\
\midrule
\textbf{This work}
  & BlueROV2 Hv.\ & Mono.\ RGB detection & \textbf{T0--T3 ablation}
  & Committed + DWA & HoloOcean & Pool, phys.\ obst., overhead GT
  & \textbf{Progressive \SZERO--\STHREE, measured}
  & \textbf{Absolute error + ranking} \\
\bottomrule
\end{tabular}
\end{table*}

